\subsection[The equivalence between graded bundles of degree $k$ and symmetric $k$-fold vector bundles]{The equivalence between graded bundles of degree $\boldsymbol{k}$\\ and symmetric $\boldsymbol{k}$-fold vector bundles}

Before giving careful statements and proofs let us sketch the basic idea. From~\cite{JG_MR_higher_vb} we know that any graded bundle $F_{k}$ can canonically be embedded in $\sT^{k}F_{k}$, and that $\sT^{k}F_{k}$ can be canonically embedded in~$\sT^{(k)}F_{k}$. \emph{Via} Proposition~\ref{prop:l_r_epi_and_inclusion} and Corollary~\ref{cor:L_emb_and_epim}, we see that the full linearisation of a~graded bundle $\Linr(F_{k})$ can be considered as a substructure of the iterated tangent bundle~$\sT^{(k)}F_{k}$. The latter has a number of canonical dif\/feomorphisms that can be enumerated by elements of the symmetric group $\Sgroup_{k}$. Thus the full linearisation of a graded bundle of degree $k$ is equipped with a number of canonical $k$-fold vector bundle morphisms $\sigma_{g}\colon \Linr(F_{k}) \rightarrow \Linr(F_{k})^{g}$, where $g \in \Sgroup_{k}$, and for any $k$-fold graded bundle $\sD =(D; \Delta_1, \ldots, \Delta_k)$ we def\/ine
\begin{gather*}\sD^g =\big(D; \Delta_{g(1)}, \ldots, \Delta_{g(k)}\big).
\end{gather*}
Note that $(\sD^{g_1})^{g_2} = \sD^{g_1g_2}$. A quick glimpse at the properties of $\sigma$'s lead us to the following def\/inition of a category~$\SymmVB[k]$, which will turn out to be equivalent with the category of graded bundles of degree $k$ (Theorem~\ref{thm:equivalence_k} below).

\begin{Definition}\label{def:SymmVB_k} Objects of the category $\SymmVB[k]$ are $k$-fold vector bundles $\sD = (D; \Delta_1, \ldots$, $\Delta_k)$ equipped with $k$-fold vector bundle isomorphisms $\sigma_g\colon \sD\rightarrow \sD^g$, where $g\in \Sgroup_k$, such that for each $g_1, g_2\in \Sgroup_k$ the composition
\begin{gather}\label{eqn:sigma1}
\sD \xrightarrow{\sigma_{g_2}} \sD^{g_2} \xrightarrow{(\sigma_{g_1})^{g_2}} \sD^{g_1g_2} = (\sD^{g_1})^{g_2}
\end{gather}
coincides with $\sigma_{g_1g_2}\colon \sD\rightarrow \sD^{g_1g_2}$ and
\begin{gather}\label{eqn:sigma2}
\sigma_{(i,j)}|_{C_{ij}} = \id_{C_{ij}},
\end{gather}
for any $1\leq i<j\leq k $, where $C_{ij}$ denotes the core bundle of the DVB $(D; \Delta_i, \Delta_j)$, and $(i,j)\in \Sgroup_k$ is the transposition swapping~$i$ and~$j$. The above $(\sigma_{g_1})^{g_2}$ coincides with $\sigma_{g_1}$ as a map $D\rightarrow D$ but we consider it here as a $k$-fold vector bundle morphism $\sD^{g_2} \rightarrow \sD^{g_1g_2}$. Note that the assignment $\sD \rightsquigarrow \sD^{g_2}$ is a functor between the categories of $k$-fold vector bundles, and it takes the mor\-phism~$\sigma_{g_1}$ to the morphism $(\sigma_{g_1})^{g_2}$. A morphism in $\SymmVB[k]$ from $(\sD, (\sigma_g)_{g\in \Sgroup_k})$ to $(\sD', ({\sigma'}_g)_{g\in \Sgroup_k})$ is a~$k$-fold vector bundle morphism $\phi\colon \sD\rightarrow \sD'$ such that
\begin{gather}\label{eqn:morphisms_in_C_k}
\sigma'_g \circ \phi = \phi^g \circ \sigma_g\colon \ \sD\rightarrow {\sD'}^g
\end{gather}
for any $g\in \Sgroup_k$.
\end{Definition}

\begin{Remark} One can give a def\/inition of symmetric $k$-fold vector bundle without setting the order of Euler vector f\/ields, and thus justifying the name ``symmetric". This is done by setting $T = \{t\colon \und{k}\to W\}$, where $\und{k}=\{1, \ldots, k\}$, $W=\{\Delta_1, \ldots, \Delta_k\}$ and considering the groupoid $\mathcal{G}\rightrightarrows T$ of all $k$-fold isomorphisms $\sD^t \to \sD^{t'}$, where $t, t'\in T$ and $\sD^t = (D; \Delta_{t(1)}, \ldots, \Delta_{t(k)})$. The symmetric group $\Sgroup_k$ acts canonically on $T$ and also on~$\mathcal{G}$. One can easily verify that giving a~groupoid morphism from the pair groupoid $T\times T\rightrightarrows T$ to $\mathcal{G}\rightrightarrows T$, covering $\id_T$, correspond with equipping $\sD$ with a family of DVB morphisms $\sigma_g$ as in Def\/inition~\ref{def:SymmVB_k} satisfying the condition~\eqref{eqn:sigma1}, but not necessary~\eqref{eqn:sigma2}.
\end{Remark}

\begin{Definition}
Let $(\sD, (\sigma_g)_{g\in \Sgroup_k})$ be a symmetric $k$-fold vector bundle, its \emph{diagonalisation} $F(\sD) \subset D$ is the graded bundle of degree $k$ constructed as the submanifold of $D$ consisting of all $\Sgroup_{k}$ invariant elements and the weight vector f\/ield is given by the total weight vector f\/ield $\Delta = \Delta_{1} + \cdots + \Delta_{k}$ restricted to the said submanifold.
\end{Definition}

It will be shortly seen that $\Delta$ is indeed tangent to $F(\sD)$ and the assignment $\sD \to F(\sD)$ is functorial. Moreover, the full linearisation functor and the diagonalisation functor are mutual inverses.

\begin{Example}
The pair $(\sT^{(k)}M, (\kappa_{g})_{g\in \Sgroup_k})$ is a symmetric $k$-fold vector bundle whose diagonalisation is clearly $\sT^{k}M$.
\end{Example}

\begin{Theorem} \label{thm:equivalence_k}
The full linearisation and the diagonalisation set up an equivalence of categories between $\GrB[k]$ and the category $\SymmVB[k]$.
\end{Theorem}
\begin{proof} Consider any graded coordinate system $\big(x^A, y^i_\veps\big)$ on $\sD$, where $1\leq i\leq N_\veps$, and the weight of $y^i_\veps$ is $\veps\in\{0,1\}^k$, $\veps\neq 0^k$. Since~$\sigma_g$ are dif\/feomorphisms of $D$ that mixes coordinates of the same total weight, the numbers $N_\veps$, $N_{\veps'}$ of coordinates of weights $\veps$ and $\veps'$ coincides if $|\veps|=|\veps'|$. Let $\Y{g}^i_\veps$ denote the same function on $D$ as $y^i_\veps$, but considered as a coordinate function on $\sD^g$, hence its weight is $\veps.g =(\veps_{g(1)}, \ldots, \veps_{g(k)})$ ($(\veps, g)\mapsto \veps.g$ is a right action of $\Sgroup_k$ on $\{0, 1\}^k$, and also on $\{0, 1\}^k \setminus \{0\}^k$). Def\/ine
\begin{gather*}%\label{eqn:modified_coord_D}
z^i_\veps := \frac{1}{k!}\sum_{g\in \Sgroup_k} \sigma^*_g \big(\Y{g}^i_{\veps.g^{-1}}\big).
\end{gather*}
Note that the weight of $\Y{g}^i_{\veps.g^{-1}}$ is $\veps.g^{-1}.g = \veps$, so $z^i_\veps$ has weight $\veps$ on $\sD$. Moreover, due to~\eqref{eqn:sigma2},
\begin{gather*}
\sigma_g^* \big(\Y{g}^i_\veps\big) = y^i_{\veps.g} + \text{lower terms}
\end{gather*}
holds for any transposition $g=(i, j)$, and so for any $g\in \Sgroup_k$, thanks to~\eqref{eqn:sigma1}, where \emph{lower terms} means the sum of products of at least two functions of non-zero weights that sum up to~$\veps.g$, hence they have to be of lower weight than $\veps.g$ with respect to the (partial) product order on~$\{0,1\}^k$.  It follows that
\begin{gather*}
z^i_\veps = y^i_\veps + \sum_{{j\geq 2, \,i_1, \ldots, i_j,}\atop {\veps^1+\cdots +\veps^j = \veps}} C^{i;\veps^1,\ldots,\veps^j}_{i_1, \ldots, i_j} y^{i_1}_{\veps^1} \cdots y^{i_j}_{\veps^j},
\end{gather*}
where the coef\/f\/icients $C^{i;\veps^1,\ldots,\veps^j}_{i_1, \ldots, i_j}$ are functions of $\big(x^A\big)$, so $\big(x^A, z^i_\veps\big)$ form a new graded coordinate system for $\sD$. Denote the corresponding coordinates on $\sD^g$ by $\Z{g}^i_\veps$. We shall show that
\begin{gather}\label{eqn:Zhi}
\sigma_h^* \big(\Z{h}^i_\veps\big) = z^i_{\veps.h}.
\end{gather}
Both sides of \eqref{eqn:Zhi} are of weight $\veps.h$, while abstracting from weights of coordinates we get
\begin{gather*}
\sigma_h^* (z^i_\veps) = \frac{1}{k!} \sum_{g\in \Sgroup_k} \sigma_h^* \sigma_g^* \big(y^i_{\veps.g^{-1}}\big) = \frac{1}{k!} \sum_{g\in \Sgroup_k} \sigma_{gh}^* \big(y^i_{\veps.g^{-1}}\big) = z^i_{\veps.h},
\end{gather*}
what justif\/ies our claim.

Thus the morphisms $\sigma_h\colon \sD\rightarrow \sD^h$ in new coordinates look extremely simple. The graded bundle $F_k\subset D$ associated with $(\sD, (\sigma_g)_{g\in \Sgroup_k})$ is given by
\begin{gather*}
F_k = \{q\in D\colon \sigma_g(q)=q \ \text{for any} \ g\in \Sgroup_k\} = \big\{\big(z^i_\veps\big)\colon z^i_\veps = z^i_{\veps'} \ \text{if}\ |\veps|=|\veps'|\big\},
\end{gather*}
and thus we may def\/ine without ambiguity graded coordinates on $F_k$ by $z^i_w := z^i_\veps|_{F_k}$, where $w=|\veps|$. Therefore, $F_k$ is indeed a graded bundle of degree $k$ and the diagonalisation is a~functor $\SymmVB[k] \to \GrB[k]$. Indeed, if $\psi\colon \sD\to \sD'$ is a morphism in~$\SymmVB[k]$, then for any $g\in\Sgroup_k$, and $q\in F_k$ holds $\sigma_g(\psi(q)) =\psi(\sigma_g(q)) = \psi(q)$ and~$\psi$ preserves the total weight given by $\Delta = \Delta_1+\cdots+\Delta_k$, hence $\psi|_{F_k}$ is a morphism in $\GrB[k]$ between diagonalisations of~$\sD$ and~$\sD'$.

Consider the full linearisation $\Linr{F_k}\subset \sT^{(k)} F_k$ with coordinates, denoted by $y^{i, (\veps)}$, inherited from $\sT^{(k)} F_k$. It is clear that $I\colon \sD\rightarrow \Linr(F_k)$, def\/ined locally by
\begin{gather*}
I^* y^{i, (\veps)} =z^i_\veps
\end{gather*}
is well def\/ined globally and makes the following diagram commutative:
\begin{gather*}
 \xymatrix{
 \sD \ar[rr]^{\sigma_g}\ar[d]_I && \sD^g \ar[d]^{I^g} \\
 \Linr F_k \ar[rr]^{\kappa_g|_{\Linr F_k}} && (\Linr F_k)^g. }
 \end{gather*}

We have described a natural one-to-one correspondence between objects in our categories. Now we shall point the correspondence of morphisms between these categories.

It is clear that a morphism of graded bundles $\phi\colon F_k\rightarrow F_k'$ gives rise to a morphism $\Linr{\phi}\colon$ $\Linr(F_k)\rightarrow \Linr(F_k')$ satisfying~\eqref{eqn:morphisms_in_C_k}. Conversely, assume $\psi\colon (\sD, (\sigma_g)) \rightarrow (\sD', (\sigma'_g))$ is a $k$-fold vector bundle morphism satisfying \eqref{eqn:morphisms_in_C_k}. It follows that we may restrict $\psi$ to $F_k\subset D$ and get a graded bundle morphism $\phi := \psi|_{F_k}\colon F_k\rightarrow F_k'$. Actually, the morphism $\psi\colon \sD\rightarrow \sD'$ is fully determined by its restriction to~$F_k$. To see this write a general form of~$\psi$ in the local coordinates $(z^i_\veps)$ def\/ined above:
\begin{gather*}
\psi^* z^{i'}_\veps = \sum_{{j\geq 1; \,k_1, \ldots, k_j,}\atop {\veps^1+\cdots +\veps^j = \veps}} Q^{i';\veps^1,\ldots,\veps^j}_{k_1, \ldots, k_j} z^{k_1}_{\veps^1} \cdots z^{k_j}_{\veps^j}, \qquad Q^{i';\veps^1,\ldots,\veps^j}_{k_1, \ldots, k_j}\in \R.
\end{gather*}
It follows from \eqref{eqn:morphisms_in_C_k} that the latter expression has to coincide with
\begin{gather*}
\sum_{{j\geq 1; \,k_1, \ldots, k_j,}\atop {\veps^1+\cdots +\veps^j = \veps}} Q^{i';\veps^1,\ldots,\veps^j}_{k_1, \ldots, k_j}   z^{k_1}_{\veps^1.g} \cdots z^{k_j}_{\veps^j.g}
\end{gather*}
for any $g\in \Sgroup_k$, so the coef\/f\/icients $Q^{\ldots}_{\ldots}$ must be symmetric in indices $k_a, k_b$ if $|\veps^a| = |\veps^b|$. Note that the formula for $\psi|_{F_k}$ involves only symmetric part of $Q^{\ldots}_{\ldots}$.
This completes our proof.
\end{proof}

Diagrammatically, we have
\begin{gather*}
\leavevmode
\begin{xy}
(0,20)*+{\catname{GrB}[k]}="a"; (50,20)*+{\catname{SymmVB}[k].}="b";%
{\ar@<1.ex>"a";"b"}?*!/_3mm/{\textnormal{full linearisation}};%
{\ar@<1.ex> "b";"a"} ?*!/_3mm/{\textnormal{diagonalisation}};%
\end{xy}
\end{gather*}

\begin{Example}
Following Example~\ref{exmpl:F3} of the full linearisation of a graded bundle of degree $3$, we see that the diagonalisation of $\Linr(F_{3})$, which is of course just $F_{3}$ itself, can be described in natural coordinates as
\begin{gather*}
 y^{a} := y^{a,(00)} = y^{a,(10)} = y^{(a),(01)}, \\ z^{i} := 2! z^{i,(10)}= 2! z^{i,(01)}= 2! z^{i,(11)}, \qquad w^{\kappa} := 3! w^{\kappa,(11)},
\end{gather*}
and by quick inspect it is clear that we recover the correct transformation laws for the non-zero weight coordinates. Up to combinatorial factors, one sets the coordinates of a given total weight equal: in a loose sense we look at the `diagonal submanifold' of the full linearisation. The generalisation of this local diagonalisation to higher degree graded bundles is clear.
\end{Example}
